\scanpage{95}
\[
[x',y']=[ax-by,bx+ay]=(a^2+b^2)[x,y]
\]
\deleted{\unsure{donc}}
\[
q(x')=q(y')=q(ax-by)=(a^2+b^2)q(x)
\]
\[
\text{donc}\qquad [x',y']/2q(x')=[x,y]/2q(x).
\]

Par suite (22) se factorise en un morphisme

\diagram{\begin{tikzcd}[ampersand replacement=\&,column sep=large,row sep=tiny]\text{(23)}\&U\arrow[r,"f"]\&\Sigma\\\&\cap\&\\\&X^{S[i]}\&\end{tikzcd}}

\uline{Proposition} \inserted{(Si $2$ inversible sur $X$)} Le morphisme précédent est un \uline{isomorphisme}.

a) C'est un monomorphisme. En effet, Supposons $\widetilde f(\unsure{\zeta|\xi})=\widetilde f(\zeta')$, où $\zeta=x+iy,\ \zeta'=x'+iy'$,
\[
\text{i.e.}\qquad [x,y]/2q(x)=[x',y']/2q(x')
\]
prouvons que $\zeta'$ est de la forme $\alpha\zeta$, $\alpha\in\Gamma G_m^{S[i]}$. Notons que $x$ et $y$ ne sont \deleted{\unsure{pas}} coll. sur aucune fibre, \inserted{\unsure{de|en} car. $\ne2$,} car si on avait $y=\lambda x$ on aurait pour $\varphi(x,y)=\lambda\varphi(x,x)=2\lambda q(x)=0$ donc $q(x)=0$, absurde
% archive sheet number: 94
